\section{Background and Motivation}
\label{sec:motivation}

\subsection{Agent Harnesses and Behavioral Contracts}

An agent harness is the non-model infrastructure that mediates between an LLM
agent and its execution
environment~\cite{langchain-harness,fowler-harness}. Harness components
include tool dispatch, memory management, orchestration, enforcement, and
feedback. This paper focuses on the enforcement and feedback components:
the mechanisms that ensure an agent's behavioral contracts hold at runtime.

Behavioral contracts are constraints that govern an agent's observable
effects rather than its internal reasoning. Projects encode them in
instruction files such as \texttt{CLAUDE.md} and
\texttt{AGENTS.md}~\cite{chatlatanagulchai2025manifests,lulla2026agentsmd}:
do not push directly to \texttt{main}; run tests before committing; never
expose secrets to the network. These contracts differ from coding-style
directives (e.g., ``prefer \texttt{const} over \texttt{let}'') in that they
produce observable effects at the system-call boundary.

\subsection{The Intent--Behavior Gap}

An agent's execution spans three abstraction levels.

\textbf{Intent} is the agent's own goals and self-constraints, i.e., what it
declares it will and will not do. 
% In ActPlane's design, the agent
% formalizes its intent in a compact DSL, just as a program formalizes its
% capability intent via \texttt{pledge()}~\cite{openbsd-pledge}. ActPlane
% does not infer intent from the agent's behavior (cf.\ AgentSight's
% passive observation~\cite{agentsight}); it receives intent as an active
% declaration that the agent itself writes and maintains.
\textbf{Action} is the tool calls the agent issues through its runtime
(\texttt{read\_file}, \texttt{run\_command}, \texttt{edit\_file}), mediated
by the agent framework.

\textbf{Behavior} is the actual OS-level operations that result:
\texttt{open}, \texttt{execve}, \texttt{connect}, and \texttt{fork}.

The mapping from action to behavior is many-to-many and opaque. A single
\texttt{run\_command("make")} triggers hundreds of system calls. The same
\texttt{git push} can be reached through a direct tool call, a
\texttt{bash~-c} string, a Python \texttt{subprocess}, or a compiled helper
binary. This opacity creates an \emph{intent--behavior gap}: contracts
declared at the intent level cannot be reliably enforced at a single
intermediate level.

Enforcement placed at the intent level (prompt instructions) is
probabilistic: LLM compliance degrades with context length and task
complexity~\cite{liu2024lost}. Enforcement placed at the action level (tool-call
guards~\cite{wang2025agentspec,shi2025progent}) is deterministic but
bypassable: it intercepts tool calls, not the OS behavior they produce.
Enforcement placed at the behavior level (syscall filters,
sandboxes~\cite{pasquier2018camquery,ciliumtetragon}) is complete but
returns opaque \texttt{-EPERM} failures with no connection to the agent's
declared contracts. Each level addresses a different part of the stack;
none spans the full gap from intent to behavior.

A cooperative agent does not intentionally bypass tool-layer guards, but
the opacity of the action-to-behavior mapping causes unintentional bypass
in routine operation. Three representative scenarios illustrate this.
(1)~An agent calls \texttt{run\_command("npm~test")}; the test harness
internally invokes \texttt{git~push} to a staging server. The tool-call
guard sees \texttt{npm~test} (a permitted action), while the kernel
observes the forbidden \texttt{git~push} behavior.
(2)~An agent calls \texttt{read\_file(".env")} to inspect a configuration
value, then later calls \texttt{run\_command("curl~api.example.com")}.
Each tool call is individually permitted, but the combination constitutes
secret exfiltration, a cross-object violation visible only to an enforcer
that tracks information flow from file read to network connect.
(3)~An agent writes a Python script containing
\texttt{subprocess.run(["git",~"push",~"--force"])} and executes it. The
tool-call guard sees \texttt{run\_command("python~script.py")}; the kernel
observes the nested \texttt{execve("git",~["push",~"--force"])}.

\subsection{Threat Model}

ActPlane targets a \emph{cooperative but probabilistic} agent. The agent is
not adversarial: it does not attempt to defeat the enforcement mechanism
through obfuscation or exploitation. However, because LLM inference is
inherently stochastic, the agent may violate any given contract with
non-zero probability at each step, regardless of context length or prompt
design. Over a long session with many steps, the cumulative probability of
at least one violation approaches certainty.

The trust boundary is the kernel. The eBPF programs and BPF maps that
propagate labels and check rules run in kernel space and cannot be modified
by the agent. The policy DSL and feedback formatting run in a trusted
userspace process (\texttt{actplane run}) that is not part of the agent's
process tree. The monitored domain is the agent's entire process tree,
including all descendants spawned through tools, shells, or subprocesses.

ActPlane does not defend against an agent that deliberately exploits kernel
vulnerabilities or manipulates the eBPF loader. Such attacks require
\texttt{CAP\_BPF} or root privileges that the agent does not hold.
Side-channel inference (deducing policy rules from \texttt{-EPERM}
patterns), deliberate label pollution (triggering over-tainting to cause
denial of service), and feedback-oracle attacks (using remediation strings
to restructure an approach toward a forbidden goal through compliant steps)
all require adversarial intent and are outside this threat model.


\section{Corpus Study}
\label{sec:corpus}

This section characterizes behavioral contracts in real agent instruction
files and identifies the subset that requires cross-object information-flow
enforcement.

\subsection{Methodology}

We collected \texttt{CLAUDE.md} and \texttt{AGENTS.md} files from 50
popular open-source projects, selected by GitHub star count and filtered
to exclude documentation-only repositories. From these files we extracted
866 candidate imperative statements using keyword matching (``never,''
``must,'' ``do not,'' ``before,'' ``only'') and manually coded each
statement along seven dimensions: (D1)~coding style vs.\ behavioral
contract, (D2)~contract category, (D3)~observable at the syscall boundary,
(D4)~expressible in the ActPlane DSL, (D5)~enforceable without false
positives under normal use, (D6)~robust to naming and path variation, and
(D7)~temptable (likely to be violated by an agent in a realistic task).
Two annotators independently coded the full set; we report Cohen's
$\kappa$ for inter-rater reliability.

\subsection{Per-Action vs.\ Cross-Object Contracts}

Behavioral contracts fall into two categories that differ in enforcement
requirements.

\textbf{Per-action contracts} constrain a single operation independently of
execution history. Examples include ``do not execute \texttt{rm~-rf}'' and
``do not run \texttt{git branch}.'' These contracts can be enforced by
matching a single system call or tool call against a static rule.

\textbf{Cross-object contracts} constrain an operation based on prior
interactions with other objects. Examples include ``a process that has read
\texttt{.env} must not connect to an external endpoint'' (data-flow
contract) and ``commit only after tests have passed'' (temporal contract).
These contracts require tracking information flow or execution ordering
across process, file, and network boundaries.

% TODO: Insert table with category breakdown and percentages once coding
% is complete.

\subsection{Implications for Enforcement}

Per-action contracts can be enforced by per-event matching at any level:
tool-call guards, seccomp filters, or eBPF kprobes. Cross-object contracts
cannot, because they depend on state accumulated across multiple operations
and objects. Enforcing them requires a mechanism that (1)~tracks provenance
across process, file, and network boundaries, (2)~encodes that provenance
as checkable state at each enforcement point, and (3)~connects violations
back to the declared contract for corrective feedback.

Labeled information-flow control satisfies all three requirements. Labels
encode cross-object provenance as per-node bitmasks, propagation rules
maintain the invariant across fork, exec, read, write, and connect
operations, and source-to-sink rules check the accumulated state at each
enforcement point. The ActPlane DSL compiles cross-object contracts
directly into this model.
